\section{Ethics, Privacy, and Child-Facing Burden}

The EEG substudy requires separate participant/caregiver information because raw neurophysiological signals are more sensitive than ordinary application telemetry and because the headband adds wearing burden and potential discomfort.

The study documentation specifies whether raw EEG is retained, pseudonymous linkage to application events, encryption, access roles, retention period, withdrawal/deletion handling, and whether sharing is public, restricted, or controlled-access.

EEG is not used to diagnose the child, infer a fixed psychological trait, or assign categorical emotions. The application remains usable without the headband. A child may request removal or discontinuation of EEG without turning that request into a behavioral failure.

The acquisition process also requires the appropriate Muse research tools and licensing pathway. The public SDK page specifically directs academic and clinical researchers to Muse Research Partners \parencite{interaxonDevelopers2026}.
